\begin{lemma}
\label{lemma-add-trivial-complex}
Suppose $R$ is a ring. Let
$$
\ldots
\xrightarrow{\varphi_{i + 1}}
R^{n_i}
\xrightarrow{\varphi_i}
R^{n_{i-1}}
\xrightarrow{\varphi_{i-1}}
\ldots
$$
be a complex of finite free $R$-modules. Suppose that for some $i$
some matrix coefficient of the map $\varphi_i$ is invertible.
Then the displayed complex is isomorphic to the direct sum of a complex
$$
\ldots \to
R^{n_{i + 2}} \xrightarrow{\varphi_{i + 2}}
R^{n_{i + 1}} \to
R^{n_i - 1} \to
R^{n_{i - 1} - 1} \to
R^{n_{i - 2}} \xrightarrow{\varphi_{i - 2}}
R^{n_{i - 3}} \to
\ldots
$$
and the complex $\ldots \to 0 \to R \to R \to 0 \to \ldots$
where the map $R \to R$ is the identity map.
\end{lemma}

\begin{proof}
The assumption means, after a change of basis of
$R^{n_i}$ and $R^{n_{i-1}}$ that the first basis
vector of $R^{n_i}$ is mapped via $\varphi_i$ to the first basis
vector of $R^{n_{i-1}}$. Let $e_j$ denote the
$j$th basis vector of $R^{n_i}$ and $f_k$ the $k$th
basis vector of $R^{n_{i-1}}$. Write $\varphi_i(e_j)
= \sum a_{jk} f_k$. So $a_{1k} = 0$ unless $k = 1$
and $a_{11} = 1$. Change basis on $R^{n_i}$ again
by setting $e'_j = e_j - a_{j1} e_1$ for $j > 1$.
After this change of coordinates we have $a_{j1} = 0$
for $j > 1$. Note the image
of $R^{n_{i + 1}} \to R^{n_i}$ is contained in the
subspace spanned by $e_j$, $j > 1$. Note also
that $R^{n_{i-1}} \to R^{n_{i-2}}$ has to annihilate
$f_1$ since it is in the image. These conditions
and the shape of the matrix $(a_{jk})$ for $\varphi_i$
imply the lemma.
\end{proof}
